\section{Die Programmiersprache Ruby}

Ruby ist eine Programmiersprache, die ab 1993 von Yukihiro Matsumoto entwickelt wurde. Dabei ließ er sich von seinen Lieblingsprogrammiersprachen Perl, Smalltalk, Eiffel, Ada und Lisp inspirieren, um eine neue Programmiersprache zu entwickeln, die sowohl funktionale und imperative Programmierung ermöglicht \citep{ruby_visual_identity_team_about_2011}. 

Eine vollständige Einführung in Ruby zu geben, würde den Rahmen dieser Diplomarbeit sprengen, weswegen ich mich auf die Herausstellung der Hauptmerkmale und Unterschiede zu anderen Sprachen konzentrieren werde, und was die Auswirkungen auf das Testgeschehen sind.

%Dem geneigten Leser seien zum Selbststudium 
% http://tryruby.org/
% .... TODO

\subsection{Einführung in Ruby}
\marginline{\includegraphics[width=0.8\marginparwidth]{material/ruby.png}
Ruby ist eine Multiparadigma-Sprache}
Ruby ist eine Multiparadigma-Sprache, die Objektorientierung, prozedurale, funktionale und nebenläufige Programmierung unterstützt. Im Gegensatz zu Java ist Ruby wie Smalltalk vollständig objektorientiert. So sind auch alle Datentypen wiederrum Objekte, was auch die primitiven Datentypen wie String und Integer umfasst. 

Ruby ist dynamisch stark typisiert, d.h. dass die Zuweisung des Typs einer Variable zur Laufzeit des Programms geschieht. Der Typ einer Variable ergibt sich damit aus ihrem Wert.

Ruby ist eine interpretierte Sprache, auch Skriptsprache genannt. Dies heisst, dass der Programmcode zur Laufzeit analysiert und ausgeführt wird. 
% TODO Quelle!!

Ruby nimmt für sich in Anspruch, eine Sprache für Menschen, und nicht für Maschinen zu sein. Dies drückt sich durch eine Syntax, die oft laut als englische Sprache gelesen werden kann. 

\setlength{\epigraphwidth}{\marginparwidth}
\marginline{\epigraph{Ruby is simple in appearance, but is very complex inside, just like our human body}{Yukihiro Matsumoto}}
\setlength{\epigraphwidth}{0.8\textwidth}

Im Nachfolgenden einige Beispiele für die Verwendung von Ruby, insbesondere die "`Alles ist ein Objekt"'-Philosophie.

\begin{lstlisting}[language=Ruby,label=Ruby Beispiele,caption=Ruby Beispiele]
 
>> puts "Hello World"
=> Hello World

>> 2.even?
=> true

>> "hallo".upcase
=> "HALLO"

>> Date.today + 2
=> #<Date: 2011-06-30>

>> a = 4 + Math.sqrt(9)
=> 7.0

>> if (0..10).include? a
>>   puts "a liegt zwischen 0 und 10"
>> end
=> a liegt zwischen 0 und 10

\end{lstlisting}

Desweiteren sei die sehr gute Unterstützung bei der Bearbeitung von Strings hervorzuheben. Insbesondere die Regulären Ausdrücke sind denen von Perl beinahe gleichmächtig und fest in die Sprache als Syntaxelement eingebaut.

\begin{lstlisting}
>> match = "2011-06-20".match /(?<year>\d{4})-\d{2}-\d{2}/
>> puts match[:year]
=> 2011
\end{lstlisting}

Auch die Bearbeitung von Arrays und listenähnlichen Strukturen ist sehr bequem, dank der Verwendungsmöglichkeit von anonymen Funktionen, bei Ruby "`Blöcke"' genannt.

\begin{lstlisting}
>> [5,5,7,3].sort
=> [3, 5, 5, 7]

# Es kann auch eine benutzerdefinierte Sortierfunktion
# angegeben werden
>> [ "string",  "rails",  "ruby" ].sort_by{ |item| item.length }
=> ["ruby", "rails", "string"]

# Die Quadratzahlen von 1 bis 5
>> (1..5).map{|element| element * 2}
=> [2, 4, 6, 8, 10]
\end{lstlisting}
Neben einem soliden objektorientierten System, bietet Ruby viele funktionale Aspekte, um die Arbeit mit Arrays und Strings sehr einfach zu gestalten. 

\paragraph{Objektsystem}
%Wichtig sei auch noch die Fähigkeit der Metaprogrammierung anzumerken.
TODO
%TODO

Im Gegensatz zu Java oder C\# dürfen Klassen zur Laufzeit um Funktionen erweitert, oder alte sogar überschrieben werden. So ist es z.B. möglich, die String-Klasse um eigene Funktionen zu erweitern.

\begin{lstlisting}
>> class String
>>   def remove_whitespace
>>     self.gsub(/\s+/, "")
>>   end
>> end

>> "Dies ist ein Test".remove_whitespace
=> "DiesisteinTest"

\end{lstlisting}

Diese Beispiele sollten als kurzer Einstieg in Ruby dienen, und einen Querschnitt durch die Besonderheiten der Sprache aufzuzeigen.

Für eine weiter Vertiefung sei das Buch "`Programming Ruby 1.9"' empfohlen, das im Detail auf die neuste Version der Programmiersprache eingeht. %TODO Buch Quelle

\subsection{Diskussion}

\paragraph{Ruby als Skriptsprache}
Dynamisch typisierte Sprachen, wie Ruby, haben gegenüber klassischen statisch typisierte Sprachen einige Nachteile. Zu allererst wir oft der Geschwindigkeitsnachteil angesprochen, den der Prozess des Interpretierens und das fehlende statische Typsystem verursachen.
Allerdings hat dies in der Regel gravierende Geschwindigkeitsnachteile. Der genaue Faktor variiert extrem, je nach Algorithmus. Ein beliebter Benchmark, shoutout.alioth, vergleicht beliebte Algorithmen der Informatik implementiert in verschiedenen Sprachen miteinander. So ergibt sich z.B. in der Gegenüberstellung von Ruby mit C ein 4-300 fache langsamere Ausführungszeit. Dem gegenüber steht allerdings jeweils nur die Hälfte bis $1/7$ der Menge an Code \citep{computer_language_benchmarks_game_ruby_2011}.

Ein Vorteil des Interpretierens, also der Übersetzung zur Laufzeit, ist eine hohe Plattformunabhängigkeit und ein leichterer Buildprozess, da das Kompilieren entfällt. 
Verfechter dynamischer Sprachen erklären weiterhin, dass diese sich ideal für prototypische Implementierungen eignen, da sich Anforderungen ständig ändern können. Weiterhin hätten Programme dynamischer Sprache eine potenziell hohe Wiederverwendbarkeit und eine höhere Lesbarkeit \citep{meijer_static_2005} \citep{ousterhout_scripting:_1998}.

Desweiteren bleiben Fehler, die der Compiler bereits entdeckt hätte, bis zur Ausführung oder schlimmstenfalls noch länger unentdeckt. Dazu gehören z.B. Tippfehler, bei denen der Wert einer nicht deklarierten Variable ausgelesen wird. Im Gegensatz zu z.B. PHP, wirft Ruby aber dann eine Exception.

Auf das Testen hat dies eine direkte Auswirkung. Viele Meinungen belegen, dass eine dynamisch typisierte Sprache mehr Tests benötigt, als eine statisch typisierte \citep{daniel_spiewak_dynamic_2010}. 



\paragraph{Metaprogrammierung}
Die Fähigkeiten zur Metaprogrammierung von Ruby bieten vielerei Möglichkeiten um Probleme effektiv zu lösen, die andernfalls nur mit erheblichem Aufwand, oder gar nicht zu lösen sind. So verwendet das beliebte Objektrelationale Datenbankframework ActiveRecord Metaprogrammierung, um einfache SQL-Statements zu erstellen.
\begin{lstlisting}
>> Person.find_by_first_name("Stefan")
  Person Load (0.2ms)  SELECT persons.* FROM persons
    WHERE users.first_name = 'Stefan' LIMIT 1
\end{lstlisting}

Die Methode \texttt{find\_by\_first\_name} existiert nicht, und wird zur Laufzeit auf Basis des Namens gebaut.

Auch das sehr beliebte Testframework Rspec, auf das wir später kurz eingehen werden, verwendet Metaprogrammierung, um Testfälle und Zusicherungen wie fast in der englischen Sprache zu formulieren. Dazu werden sämtliche Objekte von Ruby um Funktionen erweitert.

\begin{lstlisting}
>> 4.should == 3
RSpec::Expectations::ExpectationNotMetError: expected: 3
\end{lstlisting}

All diese Methoden können, richtig angewendet, zur Verbesserung der Lesbarkeit der Programme, und damit zur Erhöhrung der Wartbarkeit, führen.
Falsch angewendet stellen sie jedoch eine Gefahr da. Beim Verwenden von externen Bibliotheken, oder in einem großen Entwicklerteam können dadurch kuriose Fehler auftreten die nur sehr schwer zu finden sind. Potenziell jedes Quelltextfile kann jede Klasse zur Laufzeit verändern, ohne dass es eine Warnung gibt. 
%TODO Quelle?

\paragraph{Schlussfolgerung}

Die Verwendung von Ruby und anderen dynamischen Sprachen birgt durchaus Risiken, die zu beachten sind. Falls man diese Risiken im Kopf behält, und die Möglichkeiten der Sprache nutzt, um die Lesbarkeit zu verbessern, sind sie gerechtfertigt. Gerade bei der Entwicklung kleinerer Entwicklerteams oder Projekten mit engem Budget können dynamische Sprachen ihre Vorteile ausspielen, da sie eine eine schnellere Entwicklung ermöglicht. Im Gegensatz zu den meisten auf C basierten Sprachen, ist die Syntax von Ruby äußerst leserlich, da nur wenige Sonderzeichen verwendet werden. Auch biete Ruby mehr Funktionalität pro Programmzeile, da die Deklaration entfällt und es viel sogenannten syntaktischen Zucker gibt. Auch dies kann, richtig angewendet, der Lesbarkeit zuträglich sein.
% TODO QUelle hier was?
\epigraph{Sometimes people jot down pseudo-code on paper. If that pseudo-code runs directly on their computers, it's best, isn't it? Ruby tries to be like that, like pseudo-code that runs. }{Yukihiro Matsumoto}


\subsection{Testframeworks für Ruby}
\subsubsection{Test::Unit und Minitest}
Test::Unit (Ruby 1.8.7) und Minitest (1.9.2) sind die Testbibliotheken, die Ruby standardmäßig mitbringt. Beide basieren auf dem xUnit-, bzw. SUnit-Design von Kent Beck, und sind für Nutzer von JUnit oder NUnit leicht nachvollziehbar. %TODO

Für eine zu testende Klasse wird eine analoge Testklasse erstellt. Diese trägt per Definition denselben Namen wie die zu testende Klasse mit einem "`Test"' am Anfang. Um z.B. eine Klasse "`job"' zu testen, wird eine Datei \texttt{test\_job.rb} (Ruby Standard) oder \texttt{job\_test.rb} (Rails Standard) erstellt. Dort wiederrum wird eine Klasse mit Namen \texttt{TestJob} definiert. 

Ein Beispieltest sieht z.B. so aus:
\begin{lstlisting}
require "job"

class TestJob < Test::Unit::TestCase
  def setup
    @job = Job.create
  end
  
  def teardown
    Job.delete_all
  end
  
  def test_job_exists
    @job.title =  "Ruby on Rails Entwickler
    @job.add_location_to_title( "Dresden")
    
    assert_equal( "Ruby on Rails Entwickler in Dresden",  Job.first.title)
  end
end
\end{lstlisting}
Unsere Klasse TestJob erbt von der TestUnit Basisklasse. Sie beinhaltet die Methoden "`setup"' und "`teardown"', die jeweils vor, respektive nach jedem einzelnen Testfall aufgerufen werden.
In der Setup-Methode nehmen wir z.B. das Anlegen eines Jobs vor, in der Teardown Methode löschen wir alle Jobs in der Datenbank, um einen sauberen Test zu gewährleisten

Danach können nun beliebig viele Testmethoden folgen, deren Namen mit \texttt{test\_} beginnen müssen.
Jede Testmethode besteht in der Regel aus einer Initialisierung (kann in die setup-Methode ausgelagert werden), der Ausführung einer zu testenden Aktion und dem Prüfen der danach geltenden Eigenschaften mittels Assertions. Diese Zusicherungen sind Prädikate die oft Gleichheit oder Boolsche Rückgabewerte prüfen.

%TODO Weitere Prädikate erwähnen

Rails erledigt das Anlegen und Löschen von Testdaten selbständig. Diese werden als Fixtures bezeichnet und extern definiert. Alternativ ist der Einsatz sogenannter Factories möglich, um schnell Objekte mit bestimmten Eigenschaften zu erstellen. In jedem Fall setzt Rails die Datenbank nach jedem einzelnen Test zurück.


\subsubsection{Rspec}

